%{{{Header
\documentclass[a4paper,10pt]{article} \usepackage{anysize}
\marginsize{2cm}{2cm}{1cm}{1cm}
%\textwidth 6.0in \textheight = 664pt
\usepackage{xltxtra}
\usepackage{graphicx}
\usepackage{color}
\usepackage{xgreek}
\usepackage{fancyvrb}
\usepackage{minted}
\usepackage{listings}
\usepackage{enumitem}
\usepackage{framed}
\usepackage{relsize}
\usepackage{float}
\setmainfont[Mapping=TeX-text]{DejaVu Serif}
\newcommand{\tab}{\hspace*{3em}}

%\setmainfont{Arial}
%\setsansfont{FreeSans}
%\setmonofont{FreeMono}
\begin{document}
%\def\thesubsection {\alph{subsection}}
\renewcommand{\labelenumi}{\roman{enumi})}
\renewcommand{\labelenumii}{(\arabic{enumii})}

\begin{titlepage}
\begin{center}

\noindent\begin{minipage}[b]{.60\textwidth}
\begin{Large}
\emph{Εθνικό Μετσόβιο Πολυτεχνείο\\
Σχολή Ηλεκτρολόγων Μηχ. \& Μηχ. Υπολογιστών \\
Αλγόριθμοι και Πολυπλοκότητα\\
7\textsuperscript{ο} εξάμηνο, Ροή Λ\\
Ακαδημαϊκή περίοδος: 2010-2011\\}
\end{Large}
\end{minipage}%
\begin{minipage}[b]{.40\textwidth}
     \centering
     \includegraphics[scale=0.8]{title/ntua_logo}
\end{minipage}

\vspace{5cm}
\begin{Huge}
\underline{1\textsuperscript{η} Σειρά Γραπτών Ασκήσεων}
\end{Huge}
\vfill

\begin{flushright}
\Large{Γερακάρης Βασίλης}\\
\large{Α.Μ.: 03108092}\\
\end{flushright}
\vspace{1cm}

\large\today\\
\end{center}
\end{titlepage}


%}}}

%{{{Ταξινόμηση
\section{Ασυμπτωτικός συμβολισμός, Αναδρομικές
Σχέσεις} \setcounter{section}{1}
\begin{enumerate}

\item Ταξινόμηση
\begin{enumerate}
\item $\log {n^3}$
\item $\sqrt{n}*\log^{50}{n} $
\item $\frac{n}{\log{\log{n}} } $
\item $\log{n!} $
\item $n*\log^{10}{n} $
\item $n^{1.01} $
\item $5^{\log_{2}{n}} $
\item $\sum_{k=1}^n{k^5} = {k^6}$
\item $\log^{\log{n}}{n} = n^{\log{\log{n}} } $
\item $2^{\log_2^4{n}} $
\item $\log^{\sqrt{n}}{n} $
\item $e^{\frac{n}{\ln{n}} } $
\item $n*3^n $
\item $2^{2n} $
\item $\sqrt{n!} $
\end{enumerate}
%}}}

%{{{Τάξη μεγέθους
\item Τάξη Μεγέθους
\begin{enumerate}
\item $T(n)=5T(n/7)+n\log(n)$ ,\tab ($a=5 , b=7$) \\
$\Rightarrow n^{\log _7 {5}} = n^{0.827}$
$\Rightarrow n^{0.827} < n < n\log{n}$\\
$\Rightarrow T(n) \in \Theta(n\log{n})$

\item $T(n) = 4T(n/5)+n/\log^2{n}$ ,\tab ($a=5 , b=5$)\\
$\Rightarrow n^{\log_5 {4}} = n^{0.861}$
$\Rightarrow n^{0.861} < n/\log^2{n}$\\
$\Rightarrow T(n) \in \Theta(n/\log^2{n})$

\item $T(n) = T(n/3) +3T(n/7)+n$\\
$\Rightarrow T(n) \in \Theta(n)$

\item $T(n) = 6T(n/6)+n$\\
$\Rightarrow n^{\log _6 {6}} = n$\\
$\Rightarrow T(n) \in \Theta(n\log {n})$

\item $T(n) = T(n/3) + T(2n/3) + n$\\
$\Rightarrow T(n) \in \Theta(n\log{n})$

\item $T(n) = 16T(n/4) + n^3\log^2{n}$\\
$\Rightarrow n^{\log _4{16}} = n^2$
$\Rightarrow n^2 < n^3\log^2{n}$\\
$\Rightarrow T(n) \in \Theta(n^3\log^2{n})$

\item $T(n) = T(\sqrt{n}) + \Theta(\log{\log{n}})$ \tab
(Θέτω $k=\log{n}, f(k)=T(n))$\\
$\Rightarrow
T(\sqrt{n})=f(\sqrt{k})=f(\log{\sqrt{n}})=f(\frac{\log{n}}{2})=f(\frac{k}{2}) $\\
$\Rightarrow f(k)=f(\frac{k}{2})+\Theta(\log{k})$\\
$\Rightarrow \Theta(n)$

\item $T(n) = T(n-3) + \log{n}$\\
$\Rightarrow T(n) \in \Theta(nlogn)$

\end{enumerate}
\end{enumerate}
%}}}

%{{{Ταξινόμιση σε Πίνακα με Πολλά Ίδια Στοιχεία
\section{Ταξινόμιση σε Πίνακα με Πολλά Ίδια Στοιχεία}

%}}}

%{{{Δυαδική Αναζήτηση
\pagebreak
\section{Δυαδική Αναζήτηση}
\begin{enumerate}
\item Ξεκινάμε (για $i=1$) συγκρίνοντας το $A[i]$ με το $x$. Όσο το $x$ είναι
μεγαλύτερος από αυτό, κρατάμε την τιμή του $(i+1)$ σε μια άλλη μεταβλητή $k$ και
διπλασιάζουμε το i.
Μόλις φτάσουμε σε μεγαλύτερο αριθμό από τον $x$, εφαρμόζω ένα έλεγχο.
Αν το στοιχείο $A[i]$ είναι διάφορο του $\infty$, τότε εφαρμόζω δυαδική αναζήτηση
για το στοιχείο $x$ στο τμήμα $Α[k..i]$ του πίνακα και έχω το αποτέλεσμα για
πιθανή ύπαρξη \& θέση του στοιχείου. \\
Αν το στοιχείο $A[i]$ είναι ίσο με $\infty$, τότε ελέγχω αν $k=i$ (τερματική
συνθήκη) και αν όχι ελέγχω το $A[(k+i)/2)$. \\
\tab Αν είναι $\infty$ τότε θέτω $i= (k+i)/2$ και επαναλαβάνω.\\
\tab Αν είναι αριθμός $m$ τον συγκρίνω με το $x$.\\
\tab \tab Αν $x>m$ τότε θέτω $k= (k+i)/2$ και επαναλαμβάνω.\\
\tab \tab Αν $x<m$ τότε εφαρμόζω δυαδική αναζήτηση στο τμήμα $A[k..i]$ του
πίνακα.\\
\tab \tab Αν $x=m$ τότε είναι το ζητούμενο στοιχείο.

Μέγιστος αριθμός επαναλήψεων είναι $O(\log{n})$.

\item Παίρνω τα k πρώτα στοιχεία από κάθε πίνακα. Τα χωρίζω στη μέση και
συγκρίνω τα $A[k/2]$ με $B[k/2+1]$ και $A[k/2+1]$ με $B[k/2]$. \\
Αν $A[k/2]<B[k/2+1]$ και $B[k/2]<A[k/2+1]$ τοτε το k-οστό στοιχείο είναι το
$max\{A[k/2],B[k/2]\}$. \\
Αν $A[k/2]<B[k/2+1]$ και $B[k/2]>A[k/2+1]$ τότε επαναλαμβάνω χρησιμοποιώντας
τα στοιχεία $[A[k/2+1]..A[k]]$ και $[B[1]..B[k/2]]$.
\end{enumerate}
%}}}

%{{{Συλλογή Comics
\section{Συλλογή Comics}
Ελέχγουμε αρχικά το 1\textsuperscript{o} bit για όλα τα τεύχη, χωρίζωντάς τα
σε 2 σύνολα με πλήθος στοιχείων $\frac{n}{2}$ και $\frac{n}{2}-1$. Το τέυχος
που λείπει είναι στο σύνολο με τα λιγότερα στοιχεία, οπότε αποθηκεύουμε την
τιμή αυτού του bit. \\
Εφαρμόζουμε αναδρομικά τη μέθοδο για το επόμενο bit, κάθε φορά στο σύνολο με τα λιγότερα στοιχεία,
κάνοντας τις μισές ερωτήσεις από το προηγούμενο βήμα σε κάθε αναδρομή. \\
Τελικά, θα καταλήξουμε στον κωδικό του τεύχους που λείπει από τη συλλογή, έχοντας
ρωτήσει $n+\frac{n}{2}+\frac{n}{4}+..+1= 2n$ φορές.
%}}}

%{{{Πολυκατοικίες χωρίς Θέα
\section{Πολυκατοικίες χωρίς Θέα}
Ξεκινάμε με τον από τον $A[1]$ θέτοντας το $B[1]=0$ εφόσον δεν υπάρχει κάποιος
δυτικότερα από αυτόν. Προχωράμε στον επόμενο $A[2]$ και συγκρίνουμε το ύψος
του με τον $A[1]$. Αν $A[2]<A[1]$ τότε θέτουμε $B[2]=1$ αλλιώς συγκρίνουμε το
ύψος $A[2]$ με το ύψος του στόχου του $A[1]$. Αν πάλι δεν βρούμε ψηλότερο
τρέχουμε πάλι για το στόχο του στόχου αναδρομικά μέχρι να φτάσουμε σε 0
ή σε κάποιον ψηλότερο του $A[2]$ και τον θέτουμε ως στόχο στο πεδίο $B[2]$.
Επαναλαμβάνουμε μέχρι $n$. Μέγιστος αριθμός συγκρίσεων είναι $2n$.
%}}}

\end{document}
